% Options for packages loaded elsewhere
\PassOptionsToPackage{unicode}{hyperref}
\PassOptionsToPackage{hyphens}{url}
\documentclass[
]{article}
\usepackage{xcolor}
\usepackage[margin=1in]{geometry}
\usepackage{amsmath,amssymb}
\setcounter{secnumdepth}{-\maxdimen} % remove section numbering
\usepackage{iftex}
\ifPDFTeX
  \usepackage[T1]{fontenc}
  \usepackage[utf8]{inputenc}
  \usepackage{textcomp} % provide euro and other symbols
\else % if luatex or xetex
  \usepackage{unicode-math} % this also loads fontspec
  \defaultfontfeatures{Scale=MatchLowercase}
  \defaultfontfeatures[\rmfamily]{Ligatures=TeX,Scale=1}
\fi
\usepackage{lmodern}
\ifPDFTeX\else
  % xetex/luatex font selection
\fi
% Use upquote if available, for straight quotes in verbatim environments
\IfFileExists{upquote.sty}{\usepackage{upquote}}{}
\IfFileExists{microtype.sty}{% use microtype if available
  \usepackage[]{microtype}
  \UseMicrotypeSet[protrusion]{basicmath} % disable protrusion for tt fonts
}{}
\makeatletter
\@ifundefined{KOMAClassName}{% if non-KOMA class
  \IfFileExists{parskip.sty}{%
    \usepackage{parskip}
  }{% else
    \setlength{\parindent}{0pt}
    \setlength{\parskip}{6pt plus 2pt minus 1pt}}
}{% if KOMA class
  \KOMAoptions{parskip=half}}
\makeatother
\usepackage{color}
\usepackage{fancyvrb}
\newcommand{\VerbBar}{|}
\newcommand{\VERB}{\Verb[commandchars=\\\{\}]}
\DefineVerbatimEnvironment{Highlighting}{Verbatim}{commandchars=\\\{\}}
% Add ',fontsize=\small' for more characters per line
\usepackage{framed}
\definecolor{shadecolor}{RGB}{248,248,248}
\newenvironment{Shaded}{\begin{snugshade}}{\end{snugshade}}
\newcommand{\AlertTok}[1]{\textcolor[rgb]{0.94,0.16,0.16}{#1}}
\newcommand{\AnnotationTok}[1]{\textcolor[rgb]{0.56,0.35,0.01}{\textbf{\textit{#1}}}}
\newcommand{\AttributeTok}[1]{\textcolor[rgb]{0.13,0.29,0.53}{#1}}
\newcommand{\BaseNTok}[1]{\textcolor[rgb]{0.00,0.00,0.81}{#1}}
\newcommand{\BuiltInTok}[1]{#1}
\newcommand{\CharTok}[1]{\textcolor[rgb]{0.31,0.60,0.02}{#1}}
\newcommand{\CommentTok}[1]{\textcolor[rgb]{0.56,0.35,0.01}{\textit{#1}}}
\newcommand{\CommentVarTok}[1]{\textcolor[rgb]{0.56,0.35,0.01}{\textbf{\textit{#1}}}}
\newcommand{\ConstantTok}[1]{\textcolor[rgb]{0.56,0.35,0.01}{#1}}
\newcommand{\ControlFlowTok}[1]{\textcolor[rgb]{0.13,0.29,0.53}{\textbf{#1}}}
\newcommand{\DataTypeTok}[1]{\textcolor[rgb]{0.13,0.29,0.53}{#1}}
\newcommand{\DecValTok}[1]{\textcolor[rgb]{0.00,0.00,0.81}{#1}}
\newcommand{\DocumentationTok}[1]{\textcolor[rgb]{0.56,0.35,0.01}{\textbf{\textit{#1}}}}
\newcommand{\ErrorTok}[1]{\textcolor[rgb]{0.64,0.00,0.00}{\textbf{#1}}}
\newcommand{\ExtensionTok}[1]{#1}
\newcommand{\FloatTok}[1]{\textcolor[rgb]{0.00,0.00,0.81}{#1}}
\newcommand{\FunctionTok}[1]{\textcolor[rgb]{0.13,0.29,0.53}{\textbf{#1}}}
\newcommand{\ImportTok}[1]{#1}
\newcommand{\InformationTok}[1]{\textcolor[rgb]{0.56,0.35,0.01}{\textbf{\textit{#1}}}}
\newcommand{\KeywordTok}[1]{\textcolor[rgb]{0.13,0.29,0.53}{\textbf{#1}}}
\newcommand{\NormalTok}[1]{#1}
\newcommand{\OperatorTok}[1]{\textcolor[rgb]{0.81,0.36,0.00}{\textbf{#1}}}
\newcommand{\OtherTok}[1]{\textcolor[rgb]{0.56,0.35,0.01}{#1}}
\newcommand{\PreprocessorTok}[1]{\textcolor[rgb]{0.56,0.35,0.01}{\textit{#1}}}
\newcommand{\RegionMarkerTok}[1]{#1}
\newcommand{\SpecialCharTok}[1]{\textcolor[rgb]{0.81,0.36,0.00}{\textbf{#1}}}
\newcommand{\SpecialStringTok}[1]{\textcolor[rgb]{0.31,0.60,0.02}{#1}}
\newcommand{\StringTok}[1]{\textcolor[rgb]{0.31,0.60,0.02}{#1}}
\newcommand{\VariableTok}[1]{\textcolor[rgb]{0.00,0.00,0.00}{#1}}
\newcommand{\VerbatimStringTok}[1]{\textcolor[rgb]{0.31,0.60,0.02}{#1}}
\newcommand{\WarningTok}[1]{\textcolor[rgb]{0.56,0.35,0.01}{\textbf{\textit{#1}}}}
\usepackage{graphicx}
\makeatletter
\newsavebox\pandoc@box
\newcommand*\pandocbounded[1]{% scales image to fit in text height/width
  \sbox\pandoc@box{#1}%
  \Gscale@div\@tempa{\textheight}{\dimexpr\ht\pandoc@box+\dp\pandoc@box\relax}%
  \Gscale@div\@tempb{\linewidth}{\wd\pandoc@box}%
  \ifdim\@tempb\p@<\@tempa\p@\let\@tempa\@tempb\fi% select the smaller of both
  \ifdim\@tempa\p@<\p@\scalebox{\@tempa}{\usebox\pandoc@box}%
  \else\usebox{\pandoc@box}%
  \fi%
}
% Set default figure placement to htbp
\def\fps@figure{htbp}
\makeatother
\setlength{\emergencystretch}{3em} % prevent overfull lines
\providecommand{\tightlist}{%
  \setlength{\itemsep}{0pt}\setlength{\parskip}{0pt}}
\usepackage[spanish]{babel}
\usepackage{tikz}
\usepackage{colortbl}
\usepackage{amsmath}
\usepackage{array}
\usepackage{xcolor}
\usepackage{bookmark}
\IfFileExists{xurl.sty}{\usepackage{xurl}}{} % add URL line breaks if available
\urlstyle{same}
\hypersetup{
  hidelinks,
  pdfcreator={LaTeX via pandoc}}

\author{}
\date{\vspace{-2.5em}}

\begin{document}

\begin{Shaded}
\begin{Highlighting}[]
\FunctionTok{options}\NormalTok{(}\AttributeTok{OutDec =} \StringTok{"."}\NormalTok{)  }\CommentTok{\# Asegurar punto decimal en este chunk}

\CommentTok{\# Establecer semilla aleatoria para reproducibilidad}
\FunctionTok{set.seed}\NormalTok{(}\FunctionTok{sample}\NormalTok{(}\DecValTok{1}\SpecialCharTok{:}\DecValTok{100000}\NormalTok{, }\DecValTok{1}\NormalTok{))}

\CommentTok{\# Función para generar datos de piscinas con baldosas}
\NormalTok{generar\_datos }\OtherTok{\textless{}{-}} \ControlFlowTok{function}\NormalTok{() \{}
  \CommentTok{\# Aleatorizar el ancho de la primera piscina (múltiplo de 4 entre 4 y 12)}
\NormalTok{  ancho\_base }\OtherTok{\textless{}{-}} \FunctionTok{sample}\NormalTok{(}\FunctionTok{c}\NormalTok{(}\DecValTok{4}\NormalTok{, }\DecValTok{8}\NormalTok{, }\DecValTok{12}\NormalTok{), }\DecValTok{1}\NormalTok{)}
  
  \CommentTok{\# Generar tres piscinas con incrementos de 4 unidades}
\NormalTok{  anchos }\OtherTok{\textless{}{-}} \FunctionTok{c}\NormalTok{(ancho\_base, ancho\_base }\SpecialCharTok{+} \DecValTok{4}\NormalTok{, ancho\_base }\SpecialCharTok{+} \DecValTok{8}\NormalTok{)}
  
  \CommentTok{\# Altura fija de todas las piscinas}
\NormalTok{  altura }\OtherTok{\textless{}{-}} \DecValTok{1}
  
  \CommentTok{\# Calcular baldosas para cada piscina}
  \CommentTok{\# Fórmula: Baldosas azules = 2 * ancho + 2}
  \CommentTok{\# Baldosas blancas = 4 (siempre en las esquinas)}
\NormalTok{  baldosas\_azules }\OtherTok{\textless{}{-}} \DecValTok{2} \SpecialCharTok{*}\NormalTok{ anchos }\SpecialCharTok{+} \DecValTok{2}
\NormalTok{  baldosas\_blancas }\OtherTok{\textless{}{-}} \FunctionTok{rep}\NormalTok{(}\DecValTok{4}\NormalTok{, }\DecValTok{3}\NormalTok{)}
\NormalTok{  total\_baldosas }\OtherTok{\textless{}{-}}\NormalTok{ baldosas\_azules }\SpecialCharTok{+}\NormalTok{ baldosas\_blancas}
  
  \CommentTok{\# Aleatorizar qué piscina se pregunta (1, 2 o 3)}
\NormalTok{  piscina\_pregunta }\OtherTok{\textless{}{-}} \FunctionTok{sample}\NormalTok{(}\DecValTok{1}\SpecialCharTok{:}\DecValTok{3}\NormalTok{, }\DecValTok{1}\NormalTok{)}
  
  \CommentTok{\# Respuesta correcta: cantidad de baldosas azules de la piscina preguntada}
\NormalTok{  respuesta\_correcta }\OtherTok{\textless{}{-}}\NormalTok{ baldosas\_azules[piscina\_pregunta]}
  
  \CommentTok{\# Determinar la respuesta correcta basada en el patrón}
  \CommentTok{\# La respuesta correcta es que las baldosas azules se incrementan en 8}
  \CommentTok{\# y las blancas permanecen constantes en 4}
\NormalTok{  incremento\_azules }\OtherTok{\textless{}{-}}\NormalTok{ baldosas\_azules[}\DecValTok{2}\NormalTok{] }\SpecialCharTok{{-}}\NormalTok{ baldosas\_azules[}\DecValTok{1}\NormalTok{]}
  
  \CommentTok{\# Opciones de respuesta (afirmaciones sobre el patrón)}
  \CommentTok{\# Opción correcta: Las baldosas azules se incrementan en 8}
\NormalTok{  opcion\_correcta }\OtherTok{\textless{}{-}} \FunctionTok{paste0}\NormalTok{(}\StringTok{"Las baldosas azules se incrementan en "}\NormalTok{, incremento\_azules, }\StringTok{" de una piscina a la del siguiente tamaño."}\NormalTok{)}
  
  \CommentTok{\# Distractores (solo 3, para tener 4 opciones en total)}
\NormalTok{  distractor\_1 }\OtherTok{\textless{}{-}} \FunctionTok{paste0}\NormalTok{(}\StringTok{"Las baldosas azules se incrementan en "}\NormalTok{, incremento\_azules }\SpecialCharTok{{-}} \DecValTok{2}\NormalTok{, }\StringTok{" de una piscina a la del siguiente tamaño."}\NormalTok{)}
\NormalTok{  distractor\_2 }\OtherTok{\textless{}{-}} \StringTok{"Las baldosas blancas aumentan en ocho a medida que crece el tamaño de las piscinas."}
\NormalTok{  distractor\_3 }\OtherTok{\textless{}{-}} \StringTok{"Las baldosas azules es el doble de la cantidad de baldosas blancas en cada piscina."}
  
  \CommentTok{\# Crear vector de todas las opciones (texto) {-} solo 4 opciones}
\NormalTok{  todas\_opciones }\OtherTok{\textless{}{-}} \FunctionTok{c}\NormalTok{(opcion\_correcta, distractor\_1, distractor\_2, distractor\_3)}
  
  \CommentTok{\# Aleatorizar el orden de las opciones}
\NormalTok{  indices\_aleatorios }\OtherTok{\textless{}{-}} \FunctionTok{sample}\NormalTok{(}\DecValTok{1}\SpecialCharTok{:}\DecValTok{4}\NormalTok{, }\DecValTok{4}\NormalTok{)}
\NormalTok{  opciones\_finales }\OtherTok{\textless{}{-}}\NormalTok{ todas\_opciones[indices\_aleatorios]}
  
  \CommentTok{\# Determinar cuál es la posición correcta después de aleatorizar}
\NormalTok{  posicion\_correcta }\OtherTok{\textless{}{-}} \FunctionTok{which}\NormalTok{(indices\_aleatorios }\SpecialCharTok{==} \DecValTok{1}\NormalTok{)}
  
  \CommentTok{\# Vector de solución para r{-}exams (4 elementos)}
\NormalTok{  solucion }\OtherTok{\textless{}{-}} \FunctionTok{integer}\NormalTok{(}\DecValTok{4}\NormalTok{)}
\NormalTok{  solucion[posicion\_correcta] }\OtherTok{\textless{}{-}} \DecValTok{1}
  
  \CommentTok{\# Guardar también el incremento para la solución}
\NormalTok{  incremento\_azules\_valor }\OtherTok{\textless{}{-}}\NormalTok{ incremento\_azules}
  
  \FunctionTok{return}\NormalTok{(}\FunctionTok{list}\NormalTok{(}
    \AttributeTok{anchos =}\NormalTok{ anchos,}
    \AttributeTok{altura =}\NormalTok{ altura,}
    \AttributeTok{baldosas\_azules =}\NormalTok{ baldosas\_azules,}
    \AttributeTok{baldosas\_blancas =}\NormalTok{ baldosas\_blancas,}
    \AttributeTok{total\_baldosas =}\NormalTok{ total\_baldosas,}
    \AttributeTok{piscina\_pregunta =}\NormalTok{ piscina\_pregunta,}
    \AttributeTok{respuesta\_correcta =}\NormalTok{ respuesta\_correcta,}
    \AttributeTok{opciones\_finales =}\NormalTok{ opciones\_finales,}
    \AttributeTok{posicion\_correcta =}\NormalTok{ posicion\_correcta,}
    \AttributeTok{solucion =}\NormalTok{ solucion,}
    \AttributeTok{incremento\_azules =}\NormalTok{ incremento\_azules\_valor}
\NormalTok{  ))}
\NormalTok{\}}

\CommentTok{\# Generar datos}
\NormalTok{datos }\OtherTok{\textless{}{-}} \FunctionTok{generar\_datos}\NormalTok{()}
\end{Highlighting}
\end{Shaded}

\begin{Shaded}
\begin{Highlighting}[]
\FunctionTok{options}\NormalTok{(}\AttributeTok{OutDec =} \StringTok{"."}\NormalTok{)  }\CommentTok{\# Asegurar punto decimal en este chunk}

\CommentTok{\# Función para generar TikZ de piscinas con baldosas}
\NormalTok{generar\_tikz\_piscinas }\OtherTok{\textless{}{-}} \ControlFlowTok{function}\NormalTok{(anchos, altura, baldosas\_azules, baldosas\_blancas, piscina\_pregunta) \{}
  \CommentTok{\# Calcular el ancho máximo para centrar todas las piscinas}
\NormalTok{  ancho\_maximo }\OtherTok{\textless{}{-}} \FunctionTok{max}\NormalTok{(anchos) }\SpecialCharTok{+} \DecValTok{2}  \CommentTok{\# +2 por el borde}
  
\NormalTok{  tikz\_code }\OtherTok{\textless{}{-}} \FunctionTok{paste0}\NormalTok{(}
    \StringTok{"}\SpecialCharTok{\textbackslash{}\textbackslash{}}\StringTok{begin\{tikzpicture\}[scale=0.5]}\SpecialCharTok{\textbackslash{}n}\StringTok{"}\NormalTok{,}
    \StringTok{"\% Definir colores}\SpecialCharTok{\textbackslash{}n}\StringTok{"}\NormalTok{,}
    \StringTok{"}\SpecialCharTok{\textbackslash{}\textbackslash{}}\StringTok{definecolor\{agua\}\{RGB\}\{173,216,230\}}\SpecialCharTok{\textbackslash{}n}\StringTok{"}\NormalTok{,}
    \StringTok{"}\SpecialCharTok{\textbackslash{}\textbackslash{}}\StringTok{definecolor\{baldosa\_azul\}\{RGB\}\{0,0,255\}}\SpecialCharTok{\textbackslash{}n}\StringTok{"}\NormalTok{,}
    \StringTok{"}\SpecialCharTok{\textbackslash{}\textbackslash{}}\StringTok{definecolor\{baldosa\_blanca\}\{RGB\}\{255,255,255\}}\SpecialCharTok{\textbackslash{}n}\StringTok{"}\NormalTok{,}
    \StringTok{"}\SpecialCharTok{\textbackslash{}\textbackslash{}}\StringTok{definecolor\{borde\_negro\}\{RGB\}\{0,0,0\}}\SpecialCharTok{\textbackslash{}n\textbackslash{}n}\StringTok{"}\NormalTok{,}
    \StringTok{"\% Configuración de estilos}\SpecialCharTok{\textbackslash{}n}\StringTok{"}\NormalTok{,}
    \StringTok{"}\SpecialCharTok{\textbackslash{}\textbackslash{}}\StringTok{tikzset\{}\SpecialCharTok{\textbackslash{}n}\StringTok{"}\NormalTok{,}
    \StringTok{"  piscina/.style=\{fill=agua, draw=borde\_negro, line width=0.5pt\},}\SpecialCharTok{\textbackslash{}n}\StringTok{"}\NormalTok{,}
    \StringTok{"  baldosa\_azul/.style=\{fill=baldosa\_azul, draw=borde\_negro, line width=0.3pt\},}\SpecialCharTok{\textbackslash{}n}\StringTok{"}\NormalTok{,}
    \StringTok{"  baldosa\_blanca/.style=\{fill=baldosa\_blanca, draw=borde\_negro, line width=0.3pt\}}\SpecialCharTok{\textbackslash{}n}\StringTok{"}\NormalTok{,}
    \StringTok{"\}}\SpecialCharTok{\textbackslash{}n\textbackslash{}n}\StringTok{"}
\NormalTok{  )}
  
  \CommentTok{\# Generar las tres piscinas con mayor separación}
  \ControlFlowTok{for}\NormalTok{ (i }\ControlFlowTok{in} \DecValTok{1}\SpecialCharTok{:}\DecValTok{3}\NormalTok{) \{}
\NormalTok{    ancho }\OtherTok{\textless{}{-}}\NormalTok{ anchos[i]}
    \CommentTok{\# Espaciado vertical aumentado: cada piscina separada por 8 unidades}
\NormalTok{    y\_base }\OtherTok{\textless{}{-}}\NormalTok{ (}\DecValTok{3} \SpecialCharTok{{-}}\NormalTok{ i) }\SpecialCharTok{*} \DecValTok{8}  \CommentTok{\# Espaciado vertical aumentado}
    \CommentTok{\# Centrar horizontalmente cada piscina}
\NormalTok{    x\_offset }\OtherTok{\textless{}{-}}\NormalTok{ (ancho\_maximo }\SpecialCharTok{{-}}\NormalTok{ (ancho }\SpecialCharTok{+} \DecValTok{2}\NormalTok{)) }\SpecialCharTok{/} \DecValTok{2}
    
    \CommentTok{\# Título de la piscina}
    \ControlFlowTok{if}\NormalTok{ (i }\SpecialCharTok{==} \DecValTok{1}\NormalTok{) \{}
\NormalTok{      titulo }\OtherTok{\textless{}{-}} \StringTok{"Piscina de niños"}
\NormalTok{    \} }\ControlFlowTok{else} \ControlFlowTok{if}\NormalTok{ (i }\SpecialCharTok{==} \DecValTok{2}\NormalTok{) \{}
\NormalTok{      titulo }\OtherTok{\textless{}{-}} \StringTok{"Piscina de recreación"}
\NormalTok{    \} }\ControlFlowTok{else}\NormalTok{ \{}
\NormalTok{      titulo }\OtherTok{\textless{}{-}} \StringTok{"Piscina de entrenamiento"}
\NormalTok{    \}}
    
    \CommentTok{\# Título centrado}
\NormalTok{    x\_centro }\OtherTok{\textless{}{-}}\NormalTok{ x\_offset }\SpecialCharTok{+}\NormalTok{ (ancho }\SpecialCharTok{+} \DecValTok{2}\NormalTok{) }\SpecialCharTok{/} \DecValTok{2}
\NormalTok{    tikz\_code }\OtherTok{\textless{}{-}} \FunctionTok{paste0}\NormalTok{(tikz\_code,}
      \StringTok{"\% "}\NormalTok{, titulo, }\StringTok{"}\SpecialCharTok{\textbackslash{}n}\StringTok{"}\NormalTok{,}
      \StringTok{"}\SpecialCharTok{\textbackslash{}\textbackslash{}}\StringTok{node[above, font=}\SpecialCharTok{\textbackslash{}\textbackslash{}}\StringTok{small] at ("}\NormalTok{, x\_centro, }\StringTok{", "}\NormalTok{, y\_base }\SpecialCharTok{+}\NormalTok{ altura }\SpecialCharTok{+} \FloatTok{0.5}\NormalTok{, }\StringTok{") \{}\SpecialCharTok{\textbackslash{}\textbackslash{}}\StringTok{textbf\{Figura "}\NormalTok{, i, }\StringTok{"\}\};}\SpecialCharTok{\textbackslash{}n}\StringTok{"}\NormalTok{,}
      \StringTok{"}\SpecialCharTok{\textbackslash{}\textbackslash{}}\StringTok{node[above, font=}\SpecialCharTok{\textbackslash{}\textbackslash{}}\StringTok{scriptsize] at ("}\NormalTok{, x\_centro, }\StringTok{", "}\NormalTok{, y\_base }\SpecialCharTok{+}\NormalTok{ altura }\SpecialCharTok{+} \FloatTok{0.2}\NormalTok{, }\StringTok{") \{"}\NormalTok{, titulo, }\StringTok{"\};}\SpecialCharTok{\textbackslash{}n\textbackslash{}n}\StringTok{"}
\NormalTok{    )}
    
    \CommentTok{\# Dibujar baldosas del borde (una capa alrededor de la piscina)}
    \CommentTok{\# Ancho total con borde: ancho + 2}
    \CommentTok{\# Alto total con borde: altura + 2}
    
    \CommentTok{\# Baldosas superiores e inferiores (con offset horizontal)}
    \ControlFlowTok{for}\NormalTok{ (x }\ControlFlowTok{in} \DecValTok{0}\SpecialCharTok{:}\NormalTok{(ancho }\SpecialCharTok{+} \DecValTok{1}\NormalTok{)) \{}
\NormalTok{      x\_pos }\OtherTok{\textless{}{-}}\NormalTok{ x\_offset }\SpecialCharTok{+}\NormalTok{ x}
      \CommentTok{\# Baldosa superior}
\NormalTok{      tikz\_code }\OtherTok{\textless{}{-}} \FunctionTok{paste0}\NormalTok{(tikz\_code,}
        \StringTok{"}\SpecialCharTok{\textbackslash{}\textbackslash{}}\StringTok{draw[baldosa\_azul] ("}\NormalTok{, x\_pos, }\StringTok{", "}\NormalTok{, y\_base }\SpecialCharTok{+}\NormalTok{ altura }\SpecialCharTok{+} \DecValTok{1}\NormalTok{, }\StringTok{") rectangle ("}\NormalTok{, x\_pos }\SpecialCharTok{+} \DecValTok{1}\NormalTok{, }\StringTok{", "}\NormalTok{, y\_base }\SpecialCharTok{+}\NormalTok{ altura }\SpecialCharTok{+} \DecValTok{2}\NormalTok{, }\StringTok{");}\SpecialCharTok{\textbackslash{}n}\StringTok{"}
\NormalTok{      )}
      \CommentTok{\# Baldosa inferior}
\NormalTok{      tikz\_code }\OtherTok{\textless{}{-}} \FunctionTok{paste0}\NormalTok{(tikz\_code,}
        \StringTok{"}\SpecialCharTok{\textbackslash{}\textbackslash{}}\StringTok{draw[baldosa\_azul] ("}\NormalTok{, x\_pos, }\StringTok{", "}\NormalTok{, y\_base }\SpecialCharTok{{-}} \DecValTok{1}\NormalTok{, }\StringTok{") rectangle ("}\NormalTok{, x\_pos }\SpecialCharTok{+} \DecValTok{1}\NormalTok{, }\StringTok{", "}\NormalTok{, y\_base, }\StringTok{");}\SpecialCharTok{\textbackslash{}n}\StringTok{"}
\NormalTok{      )}
\NormalTok{    \}}
    
    \CommentTok{\# Baldosas laterales (izquierda y derecha) con offset}
    \ControlFlowTok{for}\NormalTok{ (y }\ControlFlowTok{in} \DecValTok{0}\SpecialCharTok{:}\NormalTok{(altura }\SpecialCharTok{{-}} \DecValTok{1}\NormalTok{)) \{}
      \CommentTok{\# Baldosa izquierda}
\NormalTok{      tikz\_code }\OtherTok{\textless{}{-}} \FunctionTok{paste0}\NormalTok{(tikz\_code,}
        \StringTok{"}\SpecialCharTok{\textbackslash{}\textbackslash{}}\StringTok{draw[baldosa\_azul] ("}\NormalTok{, x\_offset }\SpecialCharTok{{-}} \DecValTok{1}\NormalTok{, }\StringTok{", "}\NormalTok{, y\_base }\SpecialCharTok{+}\NormalTok{ y, }\StringTok{") rectangle ("}\NormalTok{, x\_offset, }\StringTok{", "}\NormalTok{, y\_base }\SpecialCharTok{+}\NormalTok{ y }\SpecialCharTok{+} \DecValTok{1}\NormalTok{, }\StringTok{");}\SpecialCharTok{\textbackslash{}n}\StringTok{"}
\NormalTok{      )}
      \CommentTok{\# Baldosa derecha}
\NormalTok{      tikz\_code }\OtherTok{\textless{}{-}} \FunctionTok{paste0}\NormalTok{(tikz\_code,}
        \StringTok{"}\SpecialCharTok{\textbackslash{}\textbackslash{}}\StringTok{draw[baldosa\_azul] ("}\NormalTok{, x\_offset }\SpecialCharTok{+}\NormalTok{ ancho }\SpecialCharTok{+} \DecValTok{1}\NormalTok{, }\StringTok{", "}\NormalTok{, y\_base }\SpecialCharTok{+}\NormalTok{ y, }\StringTok{") rectangle ("}\NormalTok{, x\_offset }\SpecialCharTok{+}\NormalTok{ ancho }\SpecialCharTok{+} \DecValTok{2}\NormalTok{, }\StringTok{", "}\NormalTok{, y\_base }\SpecialCharTok{+}\NormalTok{ y }\SpecialCharTok{+} \DecValTok{1}\NormalTok{, }\StringTok{");}\SpecialCharTok{\textbackslash{}n}\StringTok{"}
\NormalTok{      )}
\NormalTok{    \}}
    
    \CommentTok{\# Marcar las 4 esquinas como blancas (con offset)}
    \CommentTok{\# Esquina superior izquierda}
\NormalTok{    tikz\_code }\OtherTok{\textless{}{-}} \FunctionTok{paste0}\NormalTok{(tikz\_code,}
      \StringTok{"}\SpecialCharTok{\textbackslash{}\textbackslash{}}\StringTok{draw[baldosa\_blanca] ("}\NormalTok{, x\_offset }\SpecialCharTok{{-}} \DecValTok{1}\NormalTok{, }\StringTok{", "}\NormalTok{, y\_base }\SpecialCharTok{+}\NormalTok{ altura }\SpecialCharTok{+} \DecValTok{1}\NormalTok{, }\StringTok{") rectangle ("}\NormalTok{, x\_offset, }\StringTok{", "}\NormalTok{, y\_base }\SpecialCharTok{+}\NormalTok{ altura }\SpecialCharTok{+} \DecValTok{2}\NormalTok{, }\StringTok{");}\SpecialCharTok{\textbackslash{}n}\StringTok{"}
\NormalTok{    )}
    \CommentTok{\# Esquina superior derecha}
\NormalTok{    tikz\_code }\OtherTok{\textless{}{-}} \FunctionTok{paste0}\NormalTok{(tikz\_code,}
      \StringTok{"}\SpecialCharTok{\textbackslash{}\textbackslash{}}\StringTok{draw[baldosa\_blanca] ("}\NormalTok{, x\_offset }\SpecialCharTok{+}\NormalTok{ ancho }\SpecialCharTok{+} \DecValTok{1}\NormalTok{, }\StringTok{", "}\NormalTok{, y\_base }\SpecialCharTok{+}\NormalTok{ altura }\SpecialCharTok{+} \DecValTok{1}\NormalTok{, }\StringTok{") rectangle ("}\NormalTok{, x\_offset }\SpecialCharTok{+}\NormalTok{ ancho }\SpecialCharTok{+} \DecValTok{2}\NormalTok{, }\StringTok{", "}\NormalTok{, y\_base }\SpecialCharTok{+}\NormalTok{ altura }\SpecialCharTok{+} \DecValTok{2}\NormalTok{, }\StringTok{");}\SpecialCharTok{\textbackslash{}n}\StringTok{"}
\NormalTok{    )}
    \CommentTok{\# Esquina inferior izquierda}
\NormalTok{    tikz\_code }\OtherTok{\textless{}{-}} \FunctionTok{paste0}\NormalTok{(tikz\_code,}
      \StringTok{"}\SpecialCharTok{\textbackslash{}\textbackslash{}}\StringTok{draw[baldosa\_blanca] ("}\NormalTok{, x\_offset }\SpecialCharTok{{-}} \DecValTok{1}\NormalTok{, }\StringTok{", "}\NormalTok{, y\_base }\SpecialCharTok{{-}} \DecValTok{1}\NormalTok{, }\StringTok{") rectangle ("}\NormalTok{, x\_offset, }\StringTok{", "}\NormalTok{, y\_base, }\StringTok{");}\SpecialCharTok{\textbackslash{}n}\StringTok{"}
\NormalTok{    )}
    \CommentTok{\# Esquina inferior derecha}
\NormalTok{    tikz\_code }\OtherTok{\textless{}{-}} \FunctionTok{paste0}\NormalTok{(tikz\_code,}
      \StringTok{"}\SpecialCharTok{\textbackslash{}\textbackslash{}}\StringTok{draw[baldosa\_blanca] ("}\NormalTok{, x\_offset }\SpecialCharTok{+}\NormalTok{ ancho }\SpecialCharTok{+} \DecValTok{1}\NormalTok{, }\StringTok{", "}\NormalTok{, y\_base }\SpecialCharTok{{-}} \DecValTok{1}\NormalTok{, }\StringTok{") rectangle ("}\NormalTok{, x\_offset }\SpecialCharTok{+}\NormalTok{ ancho }\SpecialCharTok{+} \DecValTok{2}\NormalTok{, }\StringTok{", "}\NormalTok{, y\_base, }\StringTok{");}\SpecialCharTok{\textbackslash{}n}\StringTok{"}
\NormalTok{    )}
    
    \CommentTok{\# Dibujar la piscina (agua) con offset}
\NormalTok{    tikz\_code }\OtherTok{\textless{}{-}} \FunctionTok{paste0}\NormalTok{(tikz\_code,}
      \StringTok{"}\SpecialCharTok{\textbackslash{}\textbackslash{}}\StringTok{draw[piscina] ("}\NormalTok{, x\_offset, }\StringTok{", "}\NormalTok{, y\_base, }\StringTok{") rectangle ("}\NormalTok{, x\_offset }\SpecialCharTok{+}\NormalTok{ ancho }\SpecialCharTok{+} \DecValTok{1}\NormalTok{, }\StringTok{", "}\NormalTok{, y\_base }\SpecialCharTok{+}\NormalTok{ altura }\SpecialCharTok{+} \DecValTok{1}\NormalTok{, }\StringTok{");}\SpecialCharTok{\textbackslash{}n\textbackslash{}n}\StringTok{"}
\NormalTok{    )}
\NormalTok{  \}}
  
\NormalTok{  tikz\_code }\OtherTok{\textless{}{-}} \FunctionTok{paste0}\NormalTok{(tikz\_code, }\StringTok{"}\SpecialCharTok{\textbackslash{}\textbackslash{}}\StringTok{end\{tikzpicture\}"}\NormalTok{)}
  
  \FunctionTok{return}\NormalTok{(tikz\_code)}
\NormalTok{\}}

\CommentTok{\# Generar el código TikZ final}
\NormalTok{tikz\_final }\OtherTok{\textless{}{-}} \FunctionTok{generar\_tikz\_piscinas}\NormalTok{(}
\NormalTok{  datos}\SpecialCharTok{$}\NormalTok{anchos,}
\NormalTok{  datos}\SpecialCharTok{$}\NormalTok{altura,}
\NormalTok{  datos}\SpecialCharTok{$}\NormalTok{baldosas\_azules,}
\NormalTok{  datos}\SpecialCharTok{$}\NormalTok{baldosas\_blancas,}
\NormalTok{  datos}\SpecialCharTok{$}\NormalTok{piscina\_pregunta}
\NormalTok{)}
\end{Highlighting}
\end{Shaded}

\section{Question}\label{question}

Observa las siguientes figuras que muestran tres piscinas rectangulares
con diferentes tamaños. Cada piscina tiene un borde formado por baldosas
blancas (en las esquinas) y baldosas azules (en el resto del borde).

\includegraphics[width=16cm,height=\textheight,keepaspectratio]{piscinas_baldosas.png}

Basándote en el patrón observado en las figuras, ¿cuál de las siguientes
afirmaciones es correcta sobre la cantidad de baldosas?

\subsection{Answerlist}\label{answerlist}

\begin{itemize}
\tightlist
\item
  Las baldosas azules se incrementan en 6 de una piscina a la del
  siguiente tamaño.
\item
  Las baldosas azules se incrementan en 8 de una piscina a la del
  siguiente tamaño.
\item
  Las baldosas azules es el doble de la cantidad de baldosas blancas en
  cada piscina.
\item
  Las baldosas blancas aumentan en ocho a medida que crece el tamaño de
  las piscinas.
\end{itemize}

\section{Solution}\label{solution}

Para resolver este problema, debemos analizar el patrón de las baldosas
en las tres piscinas.

\textbf{Análisis del patrón:}

Observando las figuras, podemos identificar:

\begin{itemize}
\tightlist
\item
  \textbf{Figura 1 (Piscina de niños)}:

  \begin{itemize}
  \tightlist
  \item
    Ancho de la piscina: 8 unidades
  \item
    Baldosas blancas: 4 (siempre en las esquinas)
  \item
    Baldosas azules: 18
  \end{itemize}
\item
  \textbf{Figura 2 (Piscina de recreación)}:

  \begin{itemize}
  \tightlist
  \item
    Ancho de la piscina: 12 unidades
  \item
    Baldosas blancas: 4 (siempre en las esquinas)
  \item
    Baldosas azules: 26
  \end{itemize}
\item
  \textbf{Figura 3 (Piscina de entrenamiento)}:

  \begin{itemize}
  \tightlist
  \item
    Ancho de la piscina: 16 unidades
  \item
    Baldosas blancas: 4 (siempre en las esquinas)
  \item
    Baldosas azules: 34
  \end{itemize}
\end{itemize}

\textbf{Patrón identificado:}

\begin{enumerate}
\def\labelenumi{\arabic{enumi}.}
\item
  \textbf{Baldosas blancas}: Siempre son 4, independientemente del
  tamaño de la piscina (una en cada esquina del borde exterior).
\item
  \textbf{Baldosas azules}: Siguen un patrón. Si el ancho de la piscina
  es \(w\), entonces:

  \begin{itemize}
  \tightlist
  \item
    Ancho total con borde: \(w + 2\)
  \item
    Alto total con borde: \(1 + 2 = 3\)
  \item
    Total de baldosas del borde: \((w + 2) \times 3\)
  \item
    Baldosas de la piscina (agua): \(w \times 1\)
  \item
    Baldosas del borde:
    \((w + 2) \times 3 - w \times 1 = 3w + 6 - w = 2w + 6\)
  \item
    Baldosas azules: \(2w + 6 - 4 = 2w + 2\)
  \end{itemize}
\end{enumerate}

\textbf{Verificación con los datos:} - Piscina 1:
\(2 \times 8 + 2 = 18\) ✓ - Piscina 2: \(2 \times 12 + 2 = 26\) ✓ -
Piscina 3: \(2 \times 16 + 2 = 34\) ✓

\textbf{Incremento de baldosas azules:} - De la piscina 1 a la 2: 8
baldosas - De la piscina 2 a la 3: 8 baldosas

El incremento es constante: \textbf{8 baldosas azules} cada vez que el
ancho aumenta en 4 unidades.

\textbf{Evaluación de las opciones:}

Las opciones presentadas son afirmaciones sobre el patrón observado.
Analicemos cada una:

\begin{enumerate}
\def\labelenumi{\arabic{enumi}.}
\tightlist
\item
  ✅ \textbf{``Las baldosas azules se incrementan en 8''}:
  \textbf{CORRECTO}.

  \begin{itemize}
  \tightlist
  \item
    De la piscina 1 a la 2: 8 baldosas
  \item
    De la piscina 2 a la 3: 8 baldosas
  \item
    El incremento es constante e igual a 8.
  \end{itemize}
\item
  ❌ \textbf{``Las baldosas azules se incrementan en 6''}: Incorrecto.

  \begin{itemize}
  \tightlist
  \item
    El incremento real es 8, no 6.
  \end{itemize}
\item
  ❌ \textbf{``Las baldosas blancas aumentan en ocho''}: Incorrecto.

  \begin{itemize}
  \tightlist
  \item
    Las baldosas blancas siempre son 4 (una en cada esquina), no
    aumentan.
  \end{itemize}
\item
  ❌ \textbf{``Las baldosas azules es el doble de las blancas''}:
  Incorrecto.

  \begin{itemize}
  \tightlist
  \item
    Piscina 1: 18 no es el doble de 4 (8)
  \item
    Piscina 2: 26 no es el doble de 4 (8)
  \item
    Piscina 3: 34 no es el doble de 4 (8)
  \end{itemize}
\end{enumerate}

\textbf{Conclusión:} La afirmación correcta es que \textbf{las baldosas
azules se incrementan en 8} de una piscina a la del siguiente tamaño,
mientras que las baldosas blancas permanecen constantes en 4.

\subsection{Answerlist}\label{answerlist-1}

\begin{itemize}
\tightlist
\item
  Falso
\item
  Verdadero
\item
  Falso
\item
  Falso
\end{itemize}

\section{Meta-information}\label{meta-information}

exname:
piscinas\_baldosas\_patrones\_numerico\_variacional\_interpretacion\_representacion\_n2\_v1
extype: schoice exsolution: 0100 exshuffle: TRUE exsection:
Interpretación y representación exextra{[}Type{]}: Patrones y secuencias
exextra{[}Level{]}: Nivel 2 exextra{[}Language{]}: es
exextra{[}Course{]}: Matemáticas ICFES \# icfes: \# competencia:
interpretacion\_representacion \# nivel\_dificultad: 2 \# componente:
numerico\_variacional \# pensamiento: numerico \# contenido:
algebra\_generico \# eje: aplicado\_contextualizado

\end{document}
